\subsection{Vocabulary and grammar of temporal expressions}
\label{sec:generator-vocabulary}

The benchmark adopts a fixed vocabulary for writing recurrence mechanisms.
Its terminals are numerical parameters and past observations; addition,
multiplication, and the hyperbolic tangent combine them into expressions.
The symbol $t$ labels the observation being generated. Each update computes
its value from the available lag history and the mechanism's fixed parameters.
Table~\ref{tab:generator-vocabulary} defines the five elements of this language.

\begin{table}[H]
\centering\small
\begin{tabular}{@{}
>{\raggedright\arraybackslash}p{0.20\linewidth}
>{\raggedright\arraybackslash}p{0.25\linewidth}
>{\raggedright\arraybackslash}p{0.46\linewidth}@{}}
\toprule
Element & Representation & Meaning \\
\midrule
Parameter & $P_a$ &
A real coefficient supplied separately, such as $P_0=0.7$. \\
Past observation & $\operatorname{LAG}(j,L_a)$ &
Reads $x_{j,t-\ell}$ after assigning the positive integer $\ell$ to $L_a$. \\
Addition & $\operatorname{ADD}(e_1,e_2)$ &
Combines two expressions as $e_1+e_2$. \\
Multiplication & $\operatorname{MUL}(e_1,e_2)$ &
Forms $e_1e_2$: coefficients, interactions, and repeated factors. \\
Hyperbolic tangent & $\operatorname{TANH}(e)$ &
Nonlinear; output in $(-1,1)$ and derivative magnitude $\leq1$. \\
\bottomrule
\end{tabular}
\caption{Selected expression vocabulary. Parameters and lagged observations
are leaves, addition and multiplication have two children, and $\tanh$ has
one child. The source index $j$ selects the observed series.}
\label{tab:generator-vocabulary}
\end{table}

\paragraph{Composition and parameter binding.}
Let $\mathcal E$ denote any valid expression, including a complete output
tree or one of its subtrees. The recursive grammar is
\begin{equation}
\begin{aligned}
\mathcal E ::= {}& P_a \mid \operatorname{LAG}(j,L_a)\\
 &\mid\operatorname{ADD}(\mathcal E_1,\mathcal E_2)
  \mid\operatorname{MUL}(\mathcal E_1,\mathcal E_2)\\
 &\mid\operatorname{TANH}(\mathcal E_1).
\end{aligned}
\label{eq:generator-vocabulary-grammar}
\end{equation}
Here $\mathcal E_1$ and $\mathcal E_2$ are expressions built by the same rules.
The tree stores named coefficient slots $P_a$ and lag slots $L_a$; a binding
assigns real coefficients and positive integer delays. For example,
$P_0\operatorname{LAG}(0,L_0)+P_1$ becomes
$0.7x_{0,t-2}+0.01$ when $P_0=0.7$, $L_0=2$, and $P_1=0.01$.
Negative coefficients express subtraction, and repeated lag factors express
powers such as $x_{0,t-1}^2$. One output tree is retained per target series;
its lag leaves can refer to itself or another series in the system.

\paragraph{Temporal evolution through recurrence.}
Growth can arise from repeated updates of the state. For the deterministic
example
\begin{equation}
 x_t=x_{t-1}+7,\qquad t\geq1,
 \label{eq:generator-recurrent-trend}
\end{equation}
each update adds seven to the previous value. Starting from $x_0$, the
sequence is $x_0,x_0+7,x_0+14,\ldots$. Thus a trend can be represented using
only a lag leaf, a constant, and addition. This example has lag derivative
equal to one: the strict contraction conditions in
Section~\ref{sec:generator-diagnostics} select a narrower set of mechanisms.
The shared vocabulary describes expressions; each generation family adds
its own admissibility conditions.

\paragraph{Innovations, initialization, and method-specific subsets.}
The AST represents the deterministic part of the recurrence. The executor
adds the innovation $\varepsilon_{k,t}$ according to the selected noise law,
and the initial history supplies the observations required by the lags.
These choices are recorded alongside the expression. The original sampler
and bounded-domain constructor use sums and products. The Shao--Wu
construction (Section~\ref{sec:implemented-shao-wu}) wraps lag leaves in
$\tanh$ before applying its gain constraint. Heaps
(Section~\ref{sec:implemented-heaps}) expresses its linear VAR coefficients
using parameters, lag leaves, sums, and products between coefficients and
observations.

\paragraph{Alignment with the current implementation.}
The corresponding node classes and bindings are defined in
\path{src/tsfmxai/representations/expression.py}. To restrict the existing
base sampler to the selected terminals, set \texttt{time\_probability} to
zero and \texttt{require\_time} to \texttt{False}.
The current code also retains a \texttt{Time} node for earlier experiments;
the vocabulary adopted here selects the five elements in
Table~\ref{tab:generator-vocabulary}. The controlled-family comparison in
Section~\ref{sec:generator-diagnostics} already uses this restriction.
